% ESE 6180 (Fall 2026) final-project proposal.
% Two pages maximum excluding references, L4DC (PMLR) format.
\documentclass{l4dc2026}

% amsmath, amssymb, natbib, graphicx, url, algorithm2e are loaded by the class.
\usepackage{times}

\title[Evaluation-Guided Domain Randomization]{Evaluation-Guided Domain Randomization:\newline Certificate Margins and Controller Updates}

\author{%
 \Name{Shimin Chen} \Email{cshimin@engineering.upenn.edu}
}

% Blank the proceedings banner ("Proceedings of Machine Learning Research vol
% vvv:1--N, 2026") that the class puts at the top of page 1. The class applies
% page style jmlrtps there via \maketitle, so append a head-clearing step to it
% rather than emptying \jmlryear, which would also drop the copyright footer.
\makeatletter
\g@addto@macro\ps@jmlrtps{\def\@oddhead{}\let\@evenhead\@oddhead}
\makeatother

\begin{document}

\maketitle

\begin{abstract}%
Evaluation-guided domain randomization targets training conditions using a
controller's observed weaknesses, but improvement on selected conditions can
cause regression elsewhere. We propose a theory-focused study of a small
linear-system family: specialize trajectory-based certificate learning to
quadratic Lyapunov candidates, then analyze one controller update constrained
by a verified decrease margin. Cost-sensitivity bounds and counterexamples will
separate stability preservation from performance-coverage improvement. Small
numerical experiments support the analysis; cost-region expansion or a local
nonlinear extension is an ambitious additional outcome.%
\end{abstract}

\section{Introduction}

Domain randomization trains a shared controller across varying plants and
initial conditions. Evaluation can identify conditions near a performance
threshold and guide subsequent training, but proximity to that threshold need
not imply useful training signal. We study what can be certified before an
update and preserved afterward. The L4DC connection is the interaction between
statistical generalization of learned certificates and feedback-dependent
trajectory changes when the controller is updated. \\

\noindent\textbf{Statement of Contributions.} 

\section{Problem Formulation}
\label{sec:problem_formulation}

\subsection{Domain-Randomized Robust Control}
\label{subsec:general_case}
\textbf{David Example Paragraphs:} We consider the problem of domain-randomized robust control. We denote by $\theta \in \Theta$ a vector of parameters (e.g., mass, dimensions, moments of inertia) that parameterize a family of state dynamics in state $x$: 
\begin{equation}
    x_{t+1} = f(x_t, u_t; \theta),
\end{equation}
where the physical parameters are drawn from a distribution of dynamics instances $\theta \sim \mathcal{D}_\theta$. \\

The objective of domain-randomized robust control is to find a control law, $\pi$, that minimizes expected controller loss with respect to the distribution of physical systems $\mathcal{D}_\theta$. Denoting by $J$ the performance loss, we express the general case as: 
\begin{equation}
    \begin{split}
        \pi^* & := \arg\min_{\pi \in \Pi} \mathbb{E}_{\theta \sim \mathcal{D}_\theta}[J(\pi, \theta)] \\
        \text{s.t.} & \hspace{7mm} x_{t+1} = f(x_t, u_t; \theta)\\ 
        & \hspace{7mm} x_0 \sim \mathcal{D}_{0}
    \end{split}
\end{equation}

\subsection{Linear Domain-Randomized Robust Control}
\label{subsec:linear_case}
% Temporary drafting annotations. Replace the macro body by {} to hide them.
% The report's JMLR class already supplies xcolor and theorem environments.
\providecommand{\foundationsourcenote}[1]{{\normalfont\small\color{gray}%
  \hypersetup{citecolor=gray,linkcolor=gray,urlcolor=gray}%
  [Temporary source note: #1]}}

The course foundation concerns \emph{evaluation of a frozen shared gain}:
learn a quadratic candidate from trajectories, bound its violation risk on
fresh trajectories, and separately check whether it certifies all-state
stability. Controller optimization and subsequent changes to the
randomization law belong to the later stages. The following definitions
and result statements specify the foundation; proofs are deferred to
Section~\ref{sec:theory}. All vector norms are Euclidean;
$\|\cdot\|_2$ and $\|\cdot\|_F$ denote matrix operator and Frobenius norms.

\begin{definition}[Frozen evaluation experiment]
\label{def:foundation_experiment}
Let $\Theta$ be a nonempty finite set of known plants, with
$A(\theta)\in\mathbb R^{n\times n}$ and
$B(\theta)\in\mathbb R^{n\times m}$. For a shared gain
$K\in\mathbb R^{m\times n}$, define
\begin{equation}
\label{eq:foundation_closed_loop}
 x_{t+1}=A(\theta)x_t+B(\theta)u_t=G_K(\theta)x_t,
 \qquad u_t=-Kx_t,\qquad G_K(\theta)=A(\theta)-B(\theta)K.
\end{equation}
Fix $K$, positive integers $H$ (horizon) and $N$ (sample size), and a joint probability
law $\mathcal D_{\mathrm{cert}}$ on
$\Theta\times\mathcal X_0$, where
$\mathcal X_0=\{x:\|x\|_2\le r\}$ and $r>0$, before collecting data.
Draw $z_i=(\theta_i,x_{0,i})\stackrel{\mathrm{i.i.d.}}{\sim}
\mathcal D_{\mathrm{cert}}$ and record $H$ transitions of each trajectory
$\tau_i=\tau_K(z_i)$. Hold $\theta_i$ fixed within its trajectory, so
$x_{t,i}=G_K(\theta_i)^t x_{0,i}$. Set
\begin{equation}
\label{eq:foundation_dynamics_bound}
 M_K=\max_{\theta\in\Theta}\|G_K(\theta)\|_2<\infty.
\end{equation}
Dynamics and observations are exact, deterministic, full-state, and
continuous-action, with no saturation, disturbances, constraints, or early
termination. The independent sample count is $N$, not $NH$.
\foundationsourcenote{Course Chapter~3, Sections~3.1--3.2 (discrete dynamics
and feedback), and Chapter~5, Section~5.6.2 (whole trajectories as samples).
\citet{boffi2021learning}, Section~3.1, sample independent
initial states for one fixed continuous-time system. A joint plant/initial-state
law and a finite discrete horizon are our adaptations; no independence
between the two coordinates of $z_i$ is required.}
\end{definition}

\begin{definition}[Quadratic candidates and normalized decrease]
\label{def:foundation_candidates}
Fix $0<a<b<\infty$, a desired decrease coefficient $\eta>0$, and a fitting
margin $\gamma>0$, all before sampling. Define
\begin{equation}
\label{eq:foundation_candidate_class}
 \mathcal P=\{P=P^\top:aI\preceq P\preceq bI\},
 \qquad V_P(x)=x^\top Px,
 \qquad d=\frac{n(n+1)}2.
\end{equation}
The desired condition is
\begin{equation}
\label{eq:foundation_desired_decrease}
 V_P(G_K(\theta)x)-V_P(x)\le-\eta\|x\|_2^2.
\end{equation}
For a nonzero state and its observed successor define
\begin{equation}
\label{eq:foundation_residual}
 r_P(x,x^+)=\frac{x^{+\top}Px^+-x^\top Px}{\|x\|_2^2}+\eta.
\end{equation}
Then $r_P\le0$ is equivalent to~\eqref{eq:foundation_desired_decrease},
while $r_P\le-\gamma$ requires decrease at least
$(\eta+\gamma)\|x\|_2^2$. At $x=0$, the successor is zero and the desired
condition holds automatically; no normalized residual is evaluated there.
\foundationsourcenote{Course Chapter~3, Theorem~3.10, equations
(3.11)--(3.12), and Example~3.12, equations (3.20)--(3.22), supply the
discrete Lyapunov condition. Boffi Section~3.1.1 uses the unnumbered
continuous-time residual $\nabla V^\top\dot x+\alpha(V)$ in (3.1).
The actual one-step difference, state-weighted decrease, and normalization
are our choices, not an Euler substitution into that residual. The class
specializes Boffi (4.3) with feature map $m(x)=x$ and matrix $Q=P$;
the spectral bounds additionally impose positivity and fix scale.}
\end{definition}

\begin{definition}[Trajectory loss and violation risk]
\label{def:foundation_risk}
For a trajectory $\tau=(\theta,x_0,\ldots,x_H)$ containing a nonzero state
among $x_0,\ldots,x_{H-1}$, set
\begin{equation}
\label{eq:foundation_trajectory_loss}
 s_P(\tau)=\max_{\substack{0\le t<H\\x_t\ne0}}r_P(x_t,x_{t+1}),
 \qquad
 \ell_P(\tau)=\min\{1,\max\{0,1+s_P(\tau)/\gamma\}\}.
\end{equation}
For an all-zero trajectory set $s_P(\tau)=-\infty$ and $\ell_P(\tau)=0$.
Define the population loss, empirical loss, and finite-horizon violation risk
by
\begin{align}
\label{eq:foundation_population_empirical}
 L_K(P)&=\mathbb E_{z\sim\mathcal D_{\mathrm{cert}}}
                  \ell_P(\tau_K(z)),
 &\widehat L_{N,K}(P)&=\frac1N\sum_{i=1}^N\ell_P(\tau_i),\\
\label{eq:foundation_violation_risk}
 e_{K,H}(P)&=\Pr_{z\sim\mathcal D_{\mathrm{cert}}}
                   \{s_P(\tau_K(z))>0\}.
\end{align}
The risk counts a trajectory if \emph{any} of its $H$ transitions violates
the desired condition. A fresh $z$ is independent of the fitting sample.
\foundationsourcenote{Boffi (3.2) asks for a negative fitting margin;
(3.3) defines a worst-time trajectory violation probability. These become
a finite-step maximum and a joint-law probability here. Course Chapter~5,
Definitions~5.1--5.2, supply population and empirical risk; Example~5.14
motivates surrogate losses. Our exact clipped ramp replaces the margin
indicator used at the start of Boffi Section~4. Its positive loss can mean
insufficient slack rather than an actual violation.}
\end{definition}

\begin{definition}[Certificate fitting]
\label{def:foundation_fitting}
Given an optimization tolerance $\xi_{\mathrm{opt}}\ge0$, return a measurable
$\widehat P\in\mathcal P$ satisfying
\begin{equation}
\label{eq:foundation_erm}
 \widehat L_{N,K}(\widehat P)
 \le\inf_{P\in\mathcal P}\widehat L_{N,K}(P)+\xi_{\mathrm{opt}}.
\end{equation}
Only $P$ is fitted; $K$ and the recorded trajectories remain fixed.
Zero empirical loss means all nonzero sampled transitions meet the extra
margin; zero-loss feasibility is not assumed.
\foundationsourcenote{Course Chapter~5, Definition~5.3, is ERM, with
approximate optimization allowed in its discussion. This soft formulation
adapts Boffi's feasibility problem (3.2); it does not assume the feasible
zero-margin-error solution required by their Lemma~4.1. Existence of a
common quadratic certificate is not needed to state the learning problem.}
\end{definition}

\paragraph{Foundation results to establish.}
The next two lemmas specify the regularity and complexity inputs; no
stability assumption on $K$ is needed for these finite-horizon statements.

\begin{lemma}[Loss domination and parameter regularity]
\label{lem:foundation_regular_loss}
For every $P,P'\in\mathcal P$, every nonzero transition $x^+=G_K(\theta)x$,
and every trajectory of Definition~\ref{def:foundation_experiment},
\begin{align}
\label{eq:foundation_residual_lipschitz}
 |r_P(x,x^+)-r_{P'}(x,x^+)|
    &\le(M_K^2+1)\|P-P'\|_F,\\
\label{eq:foundation_loss_lipschitz}
 |\ell_P(\tau)-\ell_{P'}(\tau)|
    &\le L_{\mathrm{par}}\|P-P'\|_F,
 &L_{\mathrm{par}}&=\frac{M_K^2+1}{\gamma}.
\end{align}
Moreover $0\le\ell_P\le1$ and
$\mathbf1\{s_P(\tau)>0\}\le\ell_P(\tau)$, hence
$e_{K,H}(P)\le L_K(P)$.
\foundationsourcenote{Course Chapter~5, Proposition~5.19 and
Section~5.6.2, equation (5.42), use parameter-to-loss Lipschitz control.
Boffi Section~4's joint value/gradient norm and the unnumbered parameter
inequality in arXiv v2 Appendix~B.1 serve the same role. The constant
$(M_K^2+1)/\gamma$ and indicator domination are derived for our normalized
discrete loss, not quoted from Boffi.}
\end{lemma}

\begin{lemma}[Finite covering of the candidate class]
\label{lem:foundation_cover}
Let $\mathcal N_F(\varepsilon,\mathcal P)$ be the smallest cardinality of
a Frobenius $\varepsilon$-net with centers in $\mathcal P$.
For every $\varepsilon>0$,
\begin{equation}
\label{eq:foundation_cover}
 \|P\|_F\le b\sqrt n\quad(P\in\mathcal P),\qquad
 \log\mathcal N_F(\varepsilon,\mathcal P)
 \le d\log\!\left(1+\frac{2b\sqrt n}{\varepsilon}\right).
\end{equation}
\foundationsourcenote{Course Chapter~5, Definition~5.15 and
Propositions~5.18--5.20, define covers and bound Euclidean parameter balls.
Boffi (4.3) counts symmetric matrix parameters; v2 Appendix~B.1 uses the
unnumbered volume bound $\log\mathcal N(u;B_2^k(1))\le k\log(1+2/u)$.
Here the symmetric parameter dimension is $d=n(n+1)/2$, and packing centers
inside $\mathcal P$ preserves admissibility. This does not use the
operator-norm/dual-Sudakov refinement of their Theorem~4.3.}
\end{lemma}

\begin{theorem}[Uniform generalization for the frozen controller]
\label{thm:foundation_generalization}
For the experiment and loss defined above, fix
$\varepsilon>0$ and $\delta\in(0,1)$ before drawing the trajectories, and set
\begin{equation}
\label{eq:foundation_generalization_penalty}
 \Delta_N(\varepsilon,\delta)
 =2L_{\mathrm{par}}\varepsilon+
 \sqrt{\frac{d\log(1+2b\sqrt n/\varepsilon)+\log(2/\delta)}{2N}}.
\end{equation}
With probability at least $1-\delta$ over the $N$ independent trajectories,
simultaneously for all $P\in\mathcal P$,
\begin{equation}
\label{eq:foundation_uniform_bound}
 |L_K(P)-\widehat L_{N,K}(P)|\le\Delta_N(\varepsilon,\delta).
\end{equation}
On the same event, any fitted $\widehat P$ satisfies
\begin{align}
\label{eq:foundation_fitted_risk}
 e_{K,H}(\widehat P)
 &\le\min\{1,\widehat L_{N,K}(\widehat P)+\Delta_N(\varepsilon,\delta)\},\\
\label{eq:foundation_erm_population}
 L_K(\widehat P)
 &\le\inf_{P\in\mathcal P}L_K(P)
        +2\Delta_N(\varepsilon,\delta)+\xi_{\mathrm{opt}},
\end{align}
where the last inequality assumes~\eqref{eq:foundation_erm}.
With $\varepsilon=N^{-1/2}$ and other problem constants fixed,
$\Delta_N=O(\sqrt{(d\log(1+N)+\log(1/\delta))/N})$.
\foundationsourcenote{Course Chapter~5, Lemma~5.6, equation (5.16),
and Theorem~5.8, equation (5.20), supply Hoeffding and a finite union bound;
Proposition~5.19 transfers a parameter cover to losses. Section~5.6.2
supplies trajectory-level sampling. The approach parallels Boffi
Theorem~4.2, (4.2), but replaces Dudley's integral (v2 (B.1)) and their
fast-margin Lemma~4.1 with a two-sided concentration bound on one fixed net.
The two approximation errors give $2L_{\mathrm{par}}\varepsilon$.
This explicit elementary bound, including nonzero fitting loss and
optimization error, is our specialization; it does not reproduce Boffi's
$\widetilde O(d/N)$ or structured-matrix rate.}
\end{theorem}

The conclusion concerns violations over $H$ steps under
$\mathcal D_{\mathrm{cert}}$, not universal or infinite-horizon stability.
Its constants are independent of $H$ for this normalized homogeneous linear
construction; the trajectory loss and risk can still depend on $H$.
Changing $K$, changing the sampling law, or selecting evaluation conditions
adaptively requires fresh justified sampling or a separate transfer argument.

\begin{theorem}[Verified exponential stability]
\label{thm:foundation_verification}
Let $P\in\mathcal P$ (possibly fitted from data), and fix a nonempty
retention set $\Theta_{\mathrm{ret}}\subseteq\Theta$ before inspecting results.
For known matrices define
\begin{equation}
\label{eq:foundation_verified_margin}
 \eta_\theta(P)=\lambda_{\min}
  \!\left(P-G_K(\theta)^\top P G_K(\theta)\right),
 \qquad \eta_* =\min_{\theta\in\Theta_{\mathrm{ret}}}\eta_\theta(P).
\end{equation}
For each retained plant,~\eqref{eq:foundation_desired_decrease} holds for
every $x\in\mathbb R^n$ if and only if $\eta_\theta(P)\ge\eta$.
If $\eta_*>0$, then $0<\eta_*\le b$, every retained fixed-plant closed loop
is globally exponentially stable, and for every initial state and $t\ge1$,
\begin{equation}
\label{eq:foundation_exponential_bound}
 \|x_t\|_2^2\le\frac ba
       \left(1-\frac{\eta_*}{b}\right)^t\|x_0\|_2^2.
\end{equation}
When $\eta_*=b$, each retained loop reaches zero in one step.
Every sublevel set $\mathcal E_c(P)=\{x:x^\top Px\le c\}$, $c\ge0$,
is forward invariant for each retained plant. If also $\eta_*\ge\eta$,
the desired residual never violates on those plants at any horizon; in
particular, $e_{K,H}(P)=0$ if $\mathcal D_{\mathrm{cert}}$ is supported on
$\Theta_{\mathrm{ret}}\times\mathcal X_0$.
\foundationsourcenote{Course Chapter~3, Example~3.12, (3.22), gives
the quadratic drift identity; Theorem~3.10 and Lemma~3.11,
(3.17)--(3.19), give stability and the exponential estimate with
coefficients $a,b,\eta_*$. The eigenvalue test is our exact known-matrix
verification oracle. It is separate from Boffi (3.3)'s statistical error
and does not invoke their Theorem~5.1, (5.1)--(5.3), which instead transfers
small violation measure to a restricted decrease region under additional
incremental-stability and geometric assumptions.}
\end{theorem}

A common quadratic certificate is a restriction beyond individual plant
stability. Positive verification is conditional on the displayed test;
an unattainable fitting margin or a failed test must be reported and does
not establish instability. An initial numerical example may deliberately use
a family known to admit a common certificate. Numerical eigenvalue checks must state
their tolerance; the theorem itself is an exact-arithmetic statement.

\paragraph{Proof tasks and connection to later stages.}
Section~\ref{sec:theory} will establish the two lemmas, apply two-sided
Hoeffding concentration to a deterministic candidate net, extend to the
whole class, and specialize the quadratic Lyapunov comparison for verification.
The theorem statements already specify the constants needed by this route;
sharper complexity terms would require a separate proof and are not assumed.
Certificate fitting and verification leave $K$ unchanged. The subsequent
improvement stage will hold a verified $P$ fixed while changing $K$ and
checking how much decrease margin survives; that update is outside the
foundation's generalization theorem.

For context in those later stages, keep performance distinct from stability.
For fixed $Q,R\succeq0$, let
$J_{H,K}(\theta,x)=\sum_{t=0}^{H-1}(x_t^\top Qx_t+u_t^\top Ru_t)
=x^\top P_H(K,\theta)x$, with $x_0=x$.
Define $S_K^J=\{(\theta,x):\|x\|_2\le r,\ J_{H,K}(\theta,x)\le h\}$.
Coverage is $\mu(S_K^J)$ under a fixed reference probability measure $\mu$;
the training distribution $q$ may change. Net coverage change equals gained
mass minus lost mass. The cost matrix $P_H$ is distinct from a Lyapunov
candidate $P$, and finite-horizon acceptability does not establish stability.

\textbf{Research questions.}
What generalizes when a quadratic certificate is fitted to trajectories of a
frozen controller? Which verified margins survive an evaluation-guided update?
When can targeted training reduce its objective yet lose acceptable coverage?


% \section{Goals and Approach}

% \textbf{Stage 1: Course foundation.}
% We will adapt the certificate-learning viewpoint of \citet{boffi2021learning}
% to discrete time and $V_P(x)=x^\top Px$, with fixed
% $0<aI\preceq P\preceq bI$. For a frozen $K$, draw independent plant/initial-state
% pairs from a fixed distribution and collect $H$ transitions each. A bounded
% margin loss measures violations of
% $V_P(x_{t+1})-V_P(x_t)\le-\eta\|x_t\|_2^2$, $\eta>0$.
% A finite-cover argument will give a uniform generalization bound over the
% certificate class, treating whole trajectories as independent samples.
% This connects Chapter 5's learning tools with Chapter 3's Lyapunov theory,
% without claiming the source paper's strongest rate.
% The statistical conclusion concerns violation risk under the sampled law,
% not all-state stability. Separately, known matrices verify
% $G_K(\theta)^\top PG_K(\theta)-P\preceq-\eta_\theta I$, $\eta_\theta>0$,
% on every retained plant. Failed verification is reported.

% \textbf{Stage 2: Main theoretical investigation.}
% Hold verified $P$ fixed. For $K'=K+\Delta K$, write $G=G_K(\theta)$ and
% $d_\theta=\|B(\theta)\Delta K\|_2$. Expanding the quadratic form gives the
% sufficient preservation test
% \[
% 2\|G^\top P\|_2d_\theta+\|P\|_2d_\theta^2<\eta_\theta.
% \]
% We will connect this test to one concrete update: evaluate old costs on a fixed
% grid, mix a target near $J_{H,K}=h$ with fixed background sampling, and take an
% exact-gradient step for $F_q(K)=\mathbb E_q J_{H,K}$. Hold $q$ fixed during
% backtracking and require training descent and a positive retained margin on
% every retention plant; a capped unsuccessful search leaves $K$ unchanged.
% This preserves verified stability, not the cost-acceptable set.
% We will also derive an explicit bound
% $\|P_H(K',\theta)-P_H(K,\theta)\|_2\le L_H(\theta)\|\Delta K\|_2$,
% including the changed input penalty. With
% $b_J(\theta)=r^2L_H(\theta)\|\Delta K\|_2$, losses lie in
% $h-b_J(\theta)<J_{H,K}\le h$.
% A targeted-gradient counterexample will show that training descent and stability
% can coexist with reduced coverage. These elementary results and their limits,
% rather than a general PPO theorem, form the required analysis.

% \textbf{Stage 3: Ambitious extension.}
% Seek net expansion of $S_K^J$ under a specified update by establishing gains
% that outweigh losses. Alternatively, study a local nonlinear extension with
% an explicit remainder bound and invariant neighborhood. Choose one extension;
% Stages 1--2 remain deliverables without it. One outer update is sufficient.

% \textbf{Supporting experiments.}
% Small linear examples will compare exact and bounded drift, certificate
% verification, and gains/losses for uniform versus cost-boundary targets under
% the same update restriction. CartPole is an optional nonlinear illustration:
% with a $90^\circ$ angle cutoff and other termination rules retained, recovery
% requires 500-step survival and final-100-step RMS angle at most $5^\circ$.
% Its discrete actions and large-angle behavior do not inherit the linear theorem.
% Recovery claims require fresh episodes and simultaneous bounds
% \citep{howard2021timeuniform}. Retain raw trajectories, configurations, and seeds.

\section{Related Work}
\label{sec:related_work}

\subsection{Stability Certificates for Dynamical Systems}
\label{subsec:stability_certificate_methods}
\citet{boffi2021learning} study certificate generalization;
\citet{zhang2022adversarially} extend this to robust certificates under incremental
stability assumptions. Domain-randomized LQR \citep{fujinami2025domain}
motivates shared-gain learning. Reverse curricula \citep{florensa2017reverse}
precede performance-guided sampling, while \citet{berkenkamp2017safe} study
Lyapunov-based safe-region expansion. These results motivate our analysis but
do not imply coverage expansion for the proposed update. 


\section{Methodology}
\label{sec:method}

\section{Theoretical Results}
\label{sec:theory}

This section proves the Stage~1 results of
Section~\ref{subsec:linear_case}:
Lemmas~\ref{lem:foundation_regular_loss} and~\ref{lem:foundation_cover}
and Theorems~\ref{thm:foundation_generalization}
and~\ref{thm:foundation_verification}. The arguments are elementary and
are written to be checked line by line. Each proof opens with a short
statement of its idea. Sections~\ref{subsec:stage2_proofs}
and~\ref{subsec:stage3_envelope} add draft proofs for Stage~2 and a draft
Stage~3 starting point, both awaiting author review.

The results fit together as follows. Lemma~\ref{lem:foundation_regular_loss}
says that a small change of the certificate matrix makes only a small change
in the loss. Lemma~\ref{lem:foundation_cover} says that finitely many
matrices approximate every candidate. Theorem~\ref{thm:foundation_generalization}
combines the two with Hoeffding's inequality: concentration is proved at the
finitely many approximating matrices, and Lipschitz continuity carries it to
the whole class. Theorem~\ref{thm:foundation_verification} is independent of
the other three. It involves no data or probability and is the discrete-time
quadratic Lyapunov argument of course Chapter~3.

\subsection{Preliminaries}
\label{subsec:theory_prelims}

Let $\mathbb S^n$ denote the real symmetric $n\times n$ matrices, with
Frobenius inner product $\langle A,B\rangle_F=\operatorname{tr}(A^\top B)$
and norm $\|A\|_F=\langle A,A\rangle_F^{1/2}$. We use four standard facts.

\begin{enumerate}
\item[(F1)] \emph{Quadratic forms.} For $\Delta\in\mathbb S^n$ and
$y\in\mathbb R^n$,
\begin{equation}
\label{eq:theory_quadratic_form}
 |y^\top\Delta y|\le\|\Delta\|_F\|y\|_2^2 .
\end{equation}
Indeed, $y^\top\Delta y=\langle\Delta,yy^\top\rangle_F$, and Cauchy--Schwarz
gives $|\langle\Delta,yy^\top\rangle_F|\le\|\Delta\|_F\|yy^\top\|_F$, where
$\|yy^\top\|_F^2=\sum_{i,j}y_i^2y_j^2=\|y\|_2^4$.

\item[(F2)] \emph{Rayleigh quotient.} For $S\in\mathbb S^n$ the spectral
theorem gives an orthonormal eigenbasis, so
\begin{equation}
\label{eq:theory_rayleigh}
 x^\top Sx\ge\lambda_{\min}(S)\|x\|_2^2\quad\text{for all }x,
 \qquad\text{with equality at a unit eigenvector } v \text{ of }\lambda_{\min}(S).
\end{equation}
Applying the same fact to $-S$ gives $x^\top Sx\le\lambda_{\max}(S)\|x\|_2^2$.
Together these are the bound
$\lambda_{\min}(P)\|x\|^2\le x^\top Px\le\lambda_{\max}(P)\|x\|^2$ used in
course Chapter~3, Example~3.12.

\item[(F3)] \emph{Candidate bounds.} If $P\in\mathcal P$, then
$a\|x\|_2^2\le V_P(x)\le b\|x\|_2^2$ for every $x$. Hence every
eigenvalue of $P$ lies in $[a,b]$ (apply this at a unit eigenvector),
$\|P\|_2\le b$, and
\begin{equation}
\label{eq:theory_frobenius_bound}
 \|P\|_F^2=\operatorname{tr}(P^2)=\sum_{i=1}^n\lambda_i(P)^2\le nb^2 .
\end{equation}

\item[(F4)] \emph{Trajectories.} Within a trajectory the plant is fixed.
Hence $x_{t+1}=G_K(\theta)x_t$ and,
by~\eqref{eq:foundation_dynamics_bound},
$\|x_{t+1}\|_2\le M_K\|x_t\|_2$. In particular $x_t=0$ implies $x_s=0$ for all
$s\ge t$. The nonzero states among $x_0,\ldots,x_{H-1}$ therefore form an
initial segment $x_0,\ldots,x_{T-1}$, where
$T=T(\tau)=\#\{t<H:x_t\ne0\}\in\{0,\ldots,H\}$.
\end{enumerate}

\paragraph{Measurability.}
Give $\Theta$ its discrete $\sigma$-algebra and $\mathcal X_0$ its Borel
$\sigma$-algebra. Fix $P$. The map $z\mapsto s_P(\tau_K(z))$, with values
in $[-\infty,\infty)$, is measurable. Hence the loss $\ell_P(\tau_K(z))$
and the violation event $\{s_P(\tau_K(z))>0\}$ are measurable too, so
$L_K(P)$ and $e_{K,H}(P)$ are well defined.

To see this, first fix $\theta$ and write $G=G_K(\theta)$. The count
$T=\sum_{t<H}\mathbf 1\{G^tx_0\ne0\}$ is a sum of indicators of Borel sets,
so it is a measurable function of $x_0$. On $\{T=0\}$, $s_P=-\infty$. On
$\{T=k\}$ with $k\ge1$, $s_P$ is the maximum of the $k$ functions
$x_0\mapsto r_P(G^tx_0,G^{t+1}x_0)$, $t<k$. Each is continuous on the open
set $\{G^tx_0\ne0\}$, which contains $\{T=k\}$. So $s_P$ is Borel on each
of finitely many Borel pieces, and therefore Borel. Since $\Theta$ is
finite, the preimage of any Borel set under $z\mapsto s_P(\tau_K(z))$ is a
finite union of measurable rectangles $\{\theta\}\times(\text{Borel set})$,
which is measurable. Finally, $\ell_P=\phi(s_P)$ with the continuous $\phi$
of Step~3 below.

\subsection{Proof of Lemma~\ref{lem:foundation_regular_loss}}
\label{subsec:proof_regular_loss}

\emph{Idea.} The residual is linear in $P$, so changing $P$ changes it by a
quadratic form in $P-P'$, which (F1) controls. Taking a maximum and then
clipping are both Lipschitz operations, so the bound passes to the loss.

\begin{proof}
\emph{Step 1: residual.} Let $\Delta=P-P'\in\mathbb S^n$, let $x\ne0$, and
let $x^+=G_K(\theta)x$. The constant $\eta$ cancels in the difference:
\[
 r_P(x,x^+)-r_{P'}(x,x^+)
 =\frac{x^{+\top}\Delta x^+-x^\top\Delta x}{\|x\|_2^2}.
\]
By (F1) and $\|x^+\|_2\le M_K\|x\|_2$
(from~\eqref{eq:foundation_dynamics_bound}, since $\theta\in\Theta$), the
absolute value of the numerator is at most
\[
 \|\Delta\|_F(\|x^+\|_2^2+\|x\|_2^2)\le\|\Delta\|_F(M_K^2+1)\|x\|_2^2 .
\]
Dividing by $\|x\|_2^2$ gives~\eqref{eq:foundation_residual_lipschitz}.

\emph{Step 2: worst transition.} Fix a trajectory $\tau$. The set of times
$\{t<H:x_t\ne0\}$ is determined by $\tau$ alone and is the same for $P$ and
$P'$. If it is empty, then $\ell_P(\tau)=\ell_{P'}(\tau)=0$ by definition.
Otherwise $s_P(\tau)$ and $s_{P'}(\tau)$ are maxima over the same finite
index set. For real numbers, $|\max_t\alpha_t-\max_t\beta_t|\le
\max_t|\alpha_t-\beta_t|$. To see this, let $t^*$ attain $\max_t\alpha_t$;
then $\max_t\alpha_t-\max_t\beta_t\le\alpha_{t^*}-\beta_{t^*}$, and the other
direction is symmetric. Step~1 then gives
\[
 |s_P(\tau)-s_{P'}(\tau)|\le(M_K^2+1)\|\Delta\|_F .
\]

\emph{Step 3: clipping.} Write $\ell_P(\tau)=\phi(s_P(\tau))$, where
$\phi(s)=\min\{1,\max\{0,1+s/\gamma\}\}$ for real $s$ and
$\phi(-\infty)=0$. This matches the all-zero convention of
Definition~\ref{def:foundation_risk}. Step~2 has already handled that
case, so only real arguments matter below. The map $s\mapsto1+s/\gamma$ is
$(1/\gamma)$-Lipschitz, and clipping to $[0,1]$ is $1$-Lipschitz, so $\phi$
is $(1/\gamma)$-Lipschitz. Combining this with Step~2
gives~\eqref{eq:foundation_loss_lipschitz} with
$L_{\mathrm{par}}=(M_K^2+1)/\gamma$.

\emph{Step 4: range and domination.} The bound $0\le\ell_P\le1$ holds
because $\phi$ takes values in $[0,1]$. If $s_P(\tau)>0$, then
$1+s_P(\tau)/\gamma>1$, so $\ell_P(\tau)=1$. Otherwise the indicator is $0$,
which is at most $\ell_P(\tau)$. This covers the all-zero case
$s_P=-\infty$. So $\mathbf 1\{s_P(\tau)>0\}\le\ell_P(\tau)$ for every
$\tau$, and taking expectations under $\mathcal D_{\mathrm{cert}}$ gives
$e_{K,H}(P)\le L_K(P)$.
\end{proof}

Step~1 is the place where linear dynamics and the normalization by
$\|x\|_2^2$ are used. Because the bound is stated relative to $\|x\|_2^2$,
it does not depend on the size of the state or on the horizon $H$. This is
why Theorem~\ref{thm:foundation_generalization} has constants independent
of $H$. Step~1 holds for any symmetric $P,P'$, not only for members of
$\mathcal P$.

\subsection{Proof of Lemma~\ref{lem:foundation_cover}}
\label{subsec:proof_cover}

\emph{Idea.} We prove a volume (packing) bound in the $d$-dimensional
space $\mathbb S^n$. Course Proposition~5.18 states the analogous bound for
a Euclidean ball without proof. Course Proposition~5.20 rescales a cover of
the unit ball. Its centers are guaranteed to lie in the parameter set only
if that set contains the full ball of radius $B_\Theta$ about $0$.
Otherwise the rescaled points still cover the set, but they need not be an
$\varepsilon$-net in the sense of course Definition~5.15, which requires
centers in the set. Our $\mathcal P$ excludes a neighborhood of $0$, so it
does not contain such a ball. We therefore build the net from a maximal
separated subset of $\mathcal P$, which guarantees centers in $\mathcal P$.

\begin{proof}
\emph{Step 1: norm bound.} The first claim
in~\eqref{eq:foundation_cover} is (F3),
equation~\eqref{eq:theory_frobenius_bound}. Write $R=b\sqrt n$, so
$\mathcal P$ lies in the Frobenius ball of radius $R$ about $0$.

\emph{Step 2: coordinates.} The matrices $E_{ii}$ ($1\le i\le n$) and
$(E_{ij}+E_{ji})/\sqrt2$ ($i<j$), where $E_{ij}$ are the standard basis
matrices, form an orthonormal basis of $(\mathbb S^n,\langle\cdot,\cdot
\rangle_F)$. There are $n+n(n-1)/2=d$ of them. Expanding in this basis
identifies $\mathbb S^n$ with $\mathbb R^d$ and takes $\|\cdot\|_F$ to the
Euclidean norm. Let $\mathrm{vol}$ denote Lebesgue measure on $\mathbb R^d$.
A ball of radius $\rho$ then has volume $\rho^d\,\mathrm{vol}(B_1)$, where
$B_1$ is the unit ball.

\emph{Step 3: separated sets are small.} Call a finite set
$\{P_1,\ldots,P_k\}\subseteq\mathcal P$ \emph{$\varepsilon$-separated} if
$\|P_i-P_j\|_F>\varepsilon$ for all $i\ne j$. The open balls of radius
$\varepsilon/2$ about the $P_i$ are pairwise disjoint: a common point $Q$
would give $\|P_i-P_j\|_F\le\|P_i-Q\|_F+\|Q-P_j\|_F<\varepsilon$. Each of
these balls lies in the ball of radius $R+\varepsilon/2$ about $0$.
Comparing volumes,
\[
 k\Big(\frac\varepsilon2\Big)^d\mathrm{vol}(B_1)
 \le\Big(R+\frac\varepsilon2\Big)^d\mathrm{vol}(B_1),
 \qquad\text{so}\qquad
 k\le\Big(1+\frac{2R}{\varepsilon}\Big)^d .
\]

\emph{Step 4: a maximal separated set is a net.} The set $\mathcal P$ is
nonempty, since $aI\in\mathcal P$. By Step~3 the sizes of
$\varepsilon$-separated subsets are bounded, so we may choose one of largest
size, say $\{P_1,\ldots,P_k\}$. Every $P\in\mathcal P$ satisfies
$\|P-P_i\|_F\le\varepsilon$ for some $i$. Otherwise $P$ could be added to
the set and it would stay $\varepsilon$-separated, contradicting
maximality. Hence $\{P_i\}$ is a Frobenius $\varepsilon$-net with centers in
$\mathcal P$, and $\mathcal N_F(\varepsilon,\mathcal P)\le k\le
(1+2b\sqrt n/\varepsilon)^d$. Taking logarithms completes the proof.
\end{proof}

Steps~3--4, with $\mathcal P$ replaced by the unit ball of $\mathbb R^n$,
give course Proposition~5.18. The dimension that matters is that of the
symmetric matrices, $d=n(n+1)/2$. This matches the parameter count
$k=d(d+1)/2$ of the class (4.3) of \citet{boffi2021learning}, whose $d$ is
the feature dimension; with the feature map $m(x)=x$ their feature
dimension is our $n$, and their $k$ is our $d$. The proof
of Theorem~\ref{thm:foundation_generalization} would also work with centers
outside $\mathcal P$, since the loss is defined and bounded for every
symmetric matrix. Keeping the centers in $\mathcal P$ makes each center an
admissible certificate.

\subsection{Proof of Theorem~\ref{thm:foundation_generalization}}
\label{subsec:proof_generalization}

\emph{Idea.} This follows course Theorem~5.8 (uniform convergence over a
finite class), applied to the finite net of Lemma~\ref{lem:foundation_cover}.
Lemma~\ref{lem:foundation_regular_loss} then extends the result from the net
to all of $\mathcal P$, at a cost of $L_{\mathrm{par}}\varepsilon$ on each
of the two sides. As in course Section~5.6.2, one sample is one whole
trajectory, so the sample size is $N$ rather than $NH$.

\begin{proof}
Write $\Delta_N=\Delta_N(\varepsilon,\delta)$ and
\[
 u=\sqrt{\frac{d\log(1+2b\sqrt n/\varepsilon)+\log(2/\delta)}{2N}}>0,
 \qquad\text{so that}\qquad \Delta_N=2L_{\mathrm{par}}\varepsilon+u .
\]

\emph{Step 1: a fixed net.} Using Lemma~\ref{lem:foundation_cover}, choose
an $\varepsilon$-net $\{P^1,\ldots,P^{M_\varepsilon}\}\subseteq\mathcal P$
with $M_\varepsilon\le(1+2b\sqrt n/\varepsilon)^d$. The net depends only on $\mathcal P$ and
$\varepsilon$, which are fixed before sampling, so it does not depend on the
data. This is what allows Hoeffding's inequality to be applied at each
center (course Remark~5.7).

\emph{Step 2: concentration at one center.} Fix $j$ and set
$X_i=\ell_{P^j}(\tau_i)$, for $i=1,\ldots,N$. Each $X_i$ is a fixed
measurable function of $z_i$. The $z_i$ are i.i.d.\ draws from
$\mathcal D_{\mathrm{cert}}$, so the $X_i$ are independent, lie in $[0,1]$
by Lemma~\ref{lem:foundation_regular_loss}, and have mean $L_K(P^j)$. Their
average is $\widehat L_{N,K}(P^j)$. Course Lemma~5.6, applied once to the
$X_i$ and once to the $-X_i$, gives
\[
 \Pr\big\{|\widehat L_{N,K}(P^j)-L_K(P^j)|\ge u\big\}
 \le2\exp(-2Nu^2).
\]

\emph{Step 3: union bound.} Let $E$ be the event that
$|\widehat L_{N,K}(P^j)-L_K(P^j)|<u$ for every $j=1,\ldots,M_\varepsilon$.
As in the proof of course Theorem~5.8, but with a two-sided bound (hence
the factor $2$),
\[
 \Pr(E^c)\le2M_\varepsilon\exp(-2Nu^2)
 =2M_\varepsilon\cdot\frac{\delta}{2}
   \Big(1+\frac{2b\sqrt n}{\varepsilon}\Big)^{-d}
 \le\delta .
\]
So $\Pr(E)\ge1-\delta$. The event $E$ involves only finitely many
measurable functions of the sample, so it is measurable.

\emph{Step 4: from the net to all of $\mathcal P$.} Assume $E$ occurs, and
let $P\in\mathcal P$ be arbitrary. Choose $j$ with
$\|P-P^j\|_F\le\varepsilon$. By Lemma~\ref{lem:foundation_regular_loss},
$|\ell_P(\tau)-\ell_{P^j}(\tau)|\le L_{\mathrm{par}}\varepsilon$ for every
trajectory $\tau$. This is the pointwise form of the Lipschitz transfer in
course Proposition~5.19. Averaging the bound over the sample, or taking its
expectation, gives
\[
 |\widehat L_{N,K}(P)-\widehat L_{N,K}(P^j)|\le L_{\mathrm{par}}\varepsilon,
 \qquad
 |L_K(P)-L_K(P^j)|\le L_{\mathrm{par}}\varepsilon .
\]
The triangle inequality and the definition of $E$ then give
\[
 |L_K(P)-\widehat L_{N,K}(P)|
 \le L_{\mathrm{par}}\varepsilon+u+L_{\mathrm{par}}\varepsilon=\Delta_N .
\]
Since $P$ was arbitrary and $E$ does not depend on $P$, this
proves~\eqref{eq:foundation_uniform_bound} simultaneously for all
$P\in\mathcal P$.

\emph{Step 5: the fitted certificate.} On $E$, the bound of Step~4 holds for
every member of $\mathcal P$, including any $\widehat P$ chosen from the
data. Lemma~\ref{lem:foundation_regular_loss} gives
$e_{K,H}(\widehat P)\le L_K(\widehat P)\le\widehat L_{N,K}(\widehat P)
+\Delta_N$, and $e_{K,H}(\widehat P)\le1$ because it is a probability. This
is~\eqref{eq:foundation_fitted_risk}.

\emph{Step 6: near-optimality.} Assume also~\eqref{eq:foundation_erm}. For
any $P\in\mathcal P$, on $E$,
\begin{align*}
 L_K(\widehat P)
 &\le\widehat L_{N,K}(\widehat P)+\Delta_N
 \le\inf_{P'\in\mathcal P}\widehat L_{N,K}(P')+\xi_{\mathrm{opt}}+\Delta_N\\
 &\le\widehat L_{N,K}(P)+\xi_{\mathrm{opt}}+\Delta_N
 \le L_K(P)+2\Delta_N+\xi_{\mathrm{opt}} .
\end{align*}
The first and last inequalities use Step~4. Taking the infimum over $P$
gives~\eqref{eq:foundation_erm_population}.

\emph{Step 7: rate.} Put $\varepsilon=N^{-1/2}$ and $c=2b\sqrt n$, and let
$N\ge1$. Since $1+c\sqrt N\le(1+c)(1+N)$ and $\log(1+N)\ge\log2$,
\[
 \log(1+c\sqrt N)\le\log(1+c)+\log(1+N)
 \le\Big(1+\frac{\log(1+c)}{\log2}\Big)\log(1+N).
\]
Write $\kappa=\log(1+c)/\log2$ and
$Q_N=d\log(1+N)+\log(1/\delta)$. Since $d\ge1$,
$\log(2/\delta)\le\log(1+N)+\log(1/\delta)\le Q_N$. The numerator of $u$
is therefore at most $d(1+\kappa)\log(1+N)+Q_N\le(2+\kappa)Q_N$, so
$u\le\sqrt{1+\kappa/2}\,\sqrt{Q_N/N}$. For the first term of $\Delta_N$,
using $\log(1+N)\ge\log2$, $d\ge1$, and $\log(1/\delta)>0$,
\[
 2L_{\mathrm{par}}N^{-1/2}
 \le\frac{2L_{\mathrm{par}}}{\sqrt{\log2}}\sqrt{\frac{d\log(1+N)}{N}}
 \le\frac{2L_{\mathrm{par}}}{\sqrt{\log2}}\sqrt{\frac{Q_N}{N}} .
\]
Adding the two bounds, for every $N\ge1$ and $\delta\in(0,1)$,
\[
 \Delta_N\le C\sqrt{\frac{d\log(1+N)+\log(1/\delta)}{N}},
 \qquad
 C=\frac{2L_{\mathrm{par}}}{\sqrt{\log2}}
   +\sqrt{1+\frac{\log(1+2b\sqrt n)}{2\log2}} .
\]
The constant $C$ depends only on $b$, $n$, and $L_{\mathrm{par}}$.
\end{proof}

Two features of the proof deserve comment. First, we read ``with
probability at least $1-\delta$'' as: there is a measurable event $E$ with
$\Pr(E)\ge1-\delta$ on which all the stated bounds hold. The theorem's
phrase ``on the same event'' uses this reading. Step~4 never takes a
supremum over the uncountable class $\mathcal P$, so no measurability of a
supremum is needed, and~\eqref{eq:foundation_fitted_risk}
and~\eqref{eq:foundation_erm_population} do not need a measurable fitting
map. (The supremum in~\eqref{eq:foundation_uniform_bound} is in fact
measurable: both losses are $L_{\mathrm{par}}$-Lipschitz in $P$, so it
equals a supremum over a countable dense subset of $\mathcal P$.) The
measurability required in Definition~\ref{def:foundation_fitting} is needed
only for statements about random quantities such as
$\mathbb E[e_{K,H}(\widehat P)]$.

Second, the proof uses only that the $N$ trajectories are independent with
law $\mathcal D_{\mathrm{cert}}$, and that the net and the constants
$K,\varepsilon,\delta,\gamma,\eta$ are fixed in advance. A data-dependent
$\varepsilon$ breaks Step~1, because the net would depend on the data. A
data-dependent $K$, $\gamma$, $\eta$, or sampling law breaks Step~2,
because the $X_i$ would no longer be fixed functions of i.i.d.\ draws with
mean $L_K(P^j)$. This is why the theorem must be re-applied with fresh data
after the Stage~2 update.

\subsection{Proof of Theorem~\ref{thm:foundation_verification}}
\label{subsec:proof_verification}

\emph{Idea.} For a linear system, the change in $V_P$ along one step is the
quadratic form of $-(P-G^\top PG)$, by course Example~3.12,
equation~(3.22). That example has $Q\succ0$ from a Lyapunov equation; here
the matrix $P-G^\top PG$ plays the role of $Q$, and it is not assumed
positive definite in advance. By (F2), the smallest eigenvalue of
$P-G^\top PG$ is the exact worst-case normalized decrease. A uniform decrease, together with the
sandwich $a\|x\|^2\le V_P\le b\|x\|^2$, makes $V_P$ contract geometrically.
This is course Lemma~3.11 with $p=2$.

\begin{proof}
Fix a retained plant $\theta$. Write $G=G_K(\theta)$ and
$S_\theta=P-G^\top PG\in\mathbb S^n$, so that
$\eta_\theta(P)=\lambda_{\min}(S_\theta)$. For every $x$,
\begin{equation}
\label{eq:theory_drift_identity}
 V_P(Gx)-V_P(x)=x^\top G^\top PGx-x^\top Px=-x^\top S_\theta x .
\end{equation}

\emph{Step 1: the eigenvalue test.} By~\eqref{eq:theory_drift_identity},
condition~\eqref{eq:foundation_desired_decrease} holds for all $x$ if and
only if $x^\top S_\theta x\ge\eta\|x\|_2^2$ for all $x$. If
$\lambda_{\min}(S_\theta)\ge\eta$, this follows from (F2). Conversely, at a
unit eigenvector $v$ of $\lambda_{\min}(S_\theta)$, the inequality reads
$\lambda_{\min}(S_\theta)=v^\top S_\theta v\ge\eta$.

\emph{Step 2: $0<\eta_*\le b$.} Positivity is assumed. The minimum
in~\eqref{eq:foundation_verified_margin} is attained, since
$\Theta_{\mathrm{ret}}$ is finite and nonempty. Let $u$ be a unit
eigenvector of $\lambda_{\min}(P)$. For any retained $\theta$, (F2) and
$P\succeq0$ give
\[
 \eta_*\le\eta_\theta(P)\le u^\top S_\theta u
 =u^\top Pu-(Gu)^\top P(Gu)\le u^\top Pu=\lambda_{\min}(P)\le b .
\]
The sharper bound $\eta_*\le\lambda_{\min}(P)$ shows that the case
$\eta_*=b$ is degenerate. It forces $\lambda_{\min}(P)=b$, hence $P=bI$
because all eigenvalues lie in $[a,b]$. Then
$S_\theta=b(I-G^\top G)$ has smallest eigenvalue $b(1-\|G\|_2^2)$, which
equals $b$ only if $G_K(\theta)=0$ for every retained plant.

\emph{Step 3: geometric contraction.} For every retained $\theta$ and every
$x$, equation~\eqref{eq:theory_drift_identity} and (F2) give
$V_P(Gx)\le V_P(x)-\eta_\theta\|x\|_2^2\le V_P(x)-\eta_*\|x\|_2^2$. By (F3),
$\|x\|_2^2\ge V_P(x)/b$, so
\[
 V_P(Gx)\le\Big(1-\frac{\eta_*}{b}\Big)V_P(x),
 \qquad 0\le1-\frac{\eta_*}{b}<1,
\]
where the range of the factor comes from Step~2. Write
$\varrho=1-\eta_*/b$. Induction on $t$ gives
$V_P(x_t)\le\varrho^{\,t}V_P(x_0)$. Applying (F3) on both sides,
\[
 a\|x_t\|_2^2\le V_P(x_t)\le\varrho^{\,t}V_P(x_0)\le\varrho^{\,t}b\|x_0\|_2^2,
\]
which is~\eqref{eq:foundation_exponential_bound}. The bound also holds at
$t=0$, with the convention $\varrho^{\,0}=1$, because $b/a>1$. Taking
square roots,
$\|x_t\|_2\le\sqrt{b/a}\,(\sqrt\varrho)^{t}\|x_0\|_2$ for every initial
state. This is global exponential stability in the sense of course
Definition~3.8, with $C=\sqrt{b/a}$ and rate $\sqrt\varrho\in[0,1)$. In the
notation of course Lemma~3.11 this is the case $p=2$,
$\bar\alpha_1=a$, $\bar\alpha_2=b$, $\bar\alpha_3=\eta_*$. That lemma
needs $\bar\alpha_3>0$ (so that $\bar\alpha_3r^2$ is a class~$\mathcal K$
function), $\bar\alpha_1\le\bar\alpha_2$, and $\bar\alpha_3\le\bar\alpha_2$.
These are the hypothesis $\eta_*>0$, the assumption $a<b$, and Step~2. If $\eta_*=b$, then $\varrho=0$, so
$V_P(x_1)\le0$, and $x_1=0$ because $V_P(x_1)\ge a\|x_1\|_2^2$.

\emph{Step 4: invariant ellipsoids.} If $V_P(x)\le c$, then
$V_P(Gx)\le V_P(x)-\eta_*\|x\|_2^2\le c$. So each sublevel set
$\mathcal E_c(P)$ is mapped into itself by every retained closed loop.

\emph{Step 5: zero violation risk.} If $\eta_*\ge\eta$, then
$\eta_\theta(P)\ge\eta$ for every retained $\theta$. By Step~1, the desired
decrease holds at every state, so $r_P(x,G_K(\theta)x)\le0$ for every
$x\ne0$. Any trajectory of a retained plant therefore has $s_P(\tau)\le0$,
or $s_P(\tau)=-\infty$ if all its states are zero. Because $\Theta$ is
discrete, $\Theta_{\mathrm{ret}}\times\mathcal X_0$ is closed, so
``supported on'' means
$\mathcal D_{\mathrm{cert}}(\Theta_{\mathrm{ret}}\times\mathcal X_0)=1$. Then
$e_{K,H}(P)\le\Pr\{\theta\notin\Theta_{\mathrm{ret}}\}=0$, and this holds for
every horizon $H$.
\end{proof}

The proof uses the known matrices $A(\theta)$ and $B(\theta)$, not
sampled data. Its conclusion is therefore deterministic and does not
depend on how $P$ was obtained, including fitting by
Definition~\ref{def:foundation_fitting}. Fixing $\Theta_{\mathrm{ret}}$
before inspecting results is a protocol requirement, not a mathematical
hypothesis: the proof never uses that timing. What it protects is the
honesty of the support condition in Step~5. Since
$\mathcal D_{\mathrm{cert}}$ is a law on $\Theta\times\mathcal X_0$, that
condition is the same as requiring the $\Theta$-marginal to put mass~$1$
on $\Theta_{\mathrm{ret}}$. The fact that the decrease holds at every
state, not only on $\mathcal X_0$, would matter only if the initial-state
law were later extended beyond $\mathcal X_0$.

In floating point the theorem is applied to computed quantities, and three
checks are needed. First, the tolerance must cover the error in forming
$S_\theta=P-G^\top PG$ as well as the eigensolver error; by Weyl's
inequality the computed eigenvalue moves by at most the norm of the error
in $S_\theta$. Second, the certified lower bound $\widehat\eta_\theta-
\mathrm{tol}$, not $\widehat\eta_\theta$ itself, should be used: it must
be positive for stability, at least $\eta$ for the claim of Step~5, and it
is the value to use in the rate $\varrho$. Third, a fitted $P$ must itself
be checked to be symmetric with $aI\preceq P\preceq bI$, since (F3) and
Step~2 assume $P\in\mathcal P$ exactly.

\subsection{Draft proofs of the Stage~2 statements}
\label{subsec:stage2_proofs}

These proofs are drafts awaiting author review. Throughout, $G=G_K(\theta)$,
$E=-B(\theta)\Delta K$, and $G^+=G+E=G_{K^+}(\theta)$.

\begin{proof}[Proof of Lemma~\ref{lem:stage2_margin_drift}]
Expanding $(G+E)^\top P(G+E)-P$ gives~\eqref{eq:stage2_drift_expansion}. Let
$T=G^\top PE+E^\top PG+E^\top PE$, which is symmetric. For a unit vector $y$,
\[
 y^\top Ty=2y^\top G^\top PEy+y^\top E^\top PEy
 \le2\|G^\top P\|_2\|E\|_2+\|P\|_2\|E\|_2^2=c_\theta(\Delta K),
\]
because $\|E\|_2=d_\theta(\Delta K)$. By (F2), $T\preceq c_\theta(\Delta K)I$,
which is~\eqref{eq:stage2_drift_bound}. Since
$G^\top PG-P\preceq-\eta_\theta I$, an admissible change gives
$(G^+)^\top PG^+-P\preceq-(\eta_\theta-c_\theta(\Delta K))I
\preceq-\rho\eta_\theta I$.
\end{proof}

\begin{proof}[Proof of Theorem~\ref{thm:stage2_certified_descent}]
\emph{Accepted steps.} An accepted nonzero step has $D\ne0$, $s_j>0$, and
$\Delta K=s_jD$ admissible. Armijo~\eqref{eq:stage2_armijo} with $c_A>0$ gives
$F_q(K^+)<F_q(K)$. Lemma~\ref{lem:stage2_margin_drift} gives
\eqref{eq:stage2_preserved_margin}. Write
$\eta_*^+=\min_{\theta\in\Theta_{\mathrm{ret}}}
\lambda_{\min}(P-G_{K^+}(\theta)^\top PG_{K^+}(\theta))$ and
$\underline\eta=\rho\min_{\theta\in\Theta_{\mathrm{ret}}}\eta_\theta$, so
$\eta_*^+\ge\underline\eta>0$. Theorem~\ref{thm:foundation_verification}, applied to
$K^+$ with the same $P\in\mathcal P$, gives exponential stability, invariance of
every $\mathcal E_c(P)$, and~\eqref{eq:foundation_exponential_bound} with
$\eta_*^+$. Since $0<\underline\eta\le\eta_*^+\le b$, replacing $\eta_*^+$ by
$\underline\eta$ only enlarges the right side. If $\rho\eta_\theta\ge\eta$ on the
retention set, its eigenvalue test retains the coefficient $\eta$.

\emph{Small steps.} Let $D\ne0$ and $\varphi(s)=F_q(K+sD)$. Each $G_{K+sD}(\theta)$
is affine in $s$, so $\varphi$ is a polynomial, and
$\varphi'(0)=\langle\nabla_KF_q(K),D\rangle_F=-\|D\|_F^2$. Thus
$(\varphi(s)-\varphi(0))/s\to-\|D\|_F^2<-c_A\|D\|_F^2$ as $s\downarrow0$, since
$c_A<1$, so Armijo holds on some $(0,s_A]$. The radius holds for
$s\le\kappa/\|D\|_2$. Each $s\mapsto c_\theta(sD)$ is continuous and vanishes at
$0$, each budget $(1-\rho)\eta_\theta$ is positive, and $\Theta_{\mathrm{ret}}$ is
finite, so the consumption test holds on some $(0,s_G]$. As $s_j\to0$, an uncapped
search reaches a passing trial unless it has already accepted one. The fallback
$K^+=K$ changes nothing.
\end{proof}

\emph{Starting from the step capacity.} The midterm report and the linear studies
start the search at $\bar s=\min\{s_0,s^*(D)\}$, where $s^*(D)$ is the step capacity of
Proposition~\ref{prop:stage3_direction}, and shrink only for Armijo. The accepted-step
argument above is unchanged, because every trial is admissible and passes Armijo. The
small-step argument is also unchanged: $\bar s>0$, all smaller steps are admissible, and
Armijo holds on $(0,s_A]$. The implementation multiplies the first trial by
$1-10^{-9}$ so that rounding keeps it inside the guard.

\begin{proof}[Proof of Lemma~\ref{lem:stage2_cost_sensitivity}]
Let $\|\Delta K\|_2\le\kappa$ and $W=W_K$, $W^+=W_{K^+}$. Then
$\|G\|_2,\|G^+\|_2\le\overline M_\theta$ and $\|E\|_2\le\|B(\theta)\|_2\|\Delta K\|_2$.
Split
\begin{align*}
 P_H(K^+,\theta)-P_H(K,\theta)
 &=\sum_{t=0}^{H-1}\big((G^+)^t\big)^\top(W^+-W)(G^+)^t\\
 &\quad+\sum_{t=1}^{H-1}\Big(\big((G^+)^t\big)^\top W(G^+)^t-(G^t)^\top WG^t\Big).
\end{align*}
For the first sum, $W^+-W=K^\top R\Delta K+\Delta K^\top RK+\Delta K^\top R\Delta K$
gives $\|W^+-W\|_2\le\|R\|_2(2\|K\|_2+\kappa)\|\Delta K\|_2$. For the second, write
$A=(G^+)^t$ and $C=G^t$. Then
\[
 A^\top WA-C^\top WC=A^\top W(A-C)+(A-C)^\top WC,
 \qquad
 A-C=\sum_{j=0}^{t-1}(G^+)^{t-1-j}EG^j,
\]
so $\|A-C\|_2\le t\overline M_\theta^{\,t-1}\|E\|_2$, and the $t$-th term has norm at
most $2\|W\|_2\|B(\theta)\|_2t\overline M_\theta^{\,2t-1}\|\Delta K\|_2$. Summing
gives~\eqref{eq:stage2_cost_matrix_drift}. Then
$|x^\top(P_H(K^+,\theta)-P_H(K,\theta))x|\le r^2\mathcal L_H(\theta;\kappa)
\|\Delta K\|_2$ on $\mathcal X_0$, which is~\eqref{eq:stage2_cost_drift}.
\end{proof}

\begin{proof}[Proof of Theorem~\ref{thm:stage2_cost_regression}]
Write $b=b_J(\theta;\Delta K)$. If $J_{H,K}(\theta,x)\le h-b$,
then~\eqref{eq:stage2_cost_drift} gives $J_{H,K^+}(\theta,x)\le h$,
proving~\eqref{eq:stage2_interior_preservation}. If $z$ is lost, then
$J_{H,K}(z)\le h<J_{H,K^+}(z)\le J_{H,K}(z)+b$, so $z$ lies in the loss band,
proving~\eqref{eq:stage2_losses}. Costs are continuous in $x$ and $\Theta$ is finite,
so these sets are measurable. Finally
$\mu(S_{K^+}^J)=\mu(S_{K^+}^J\cap S_K^J)+\mu(S_{K^+}^J\setminus S_K^J)$, likewise
for $S_K^J$; subtracting gives the identity
in~\eqref{eq:stage2_coverage_accounting}, and dropping the gained mass
with~\eqref{eq:stage2_losses} gives the inequality.
\end{proof}

\begin{proof}[Proof of Proposition~\ref{prop:stage2_counterexample}]
Here $G_{\widetilde K}(\theta)=\theta-\widetilde K$ and
$J_{2,\widetilde K}(\theta,1)=1+(\theta-\widetilde K)^2$. At $K=1/2$ both costs are
$5/4=h$, so $C_J(K)=1$. With $q=(1/4,3/4)$,
$F_q(\widetilde K)=\tfrac14(1+\widetilde K^2)+\tfrac34(1+(1-\widetilde K)^2)$, so
$F_q(1/2)=1.25$, $\nabla F_q(1/2)=1/4-3/4=-1/2$, and $D=1/2$. Then $s_0=2/5$ gives
$\Delta K=1/5\le\kappa$ and $K^+=7/10$. Both old margins are $1-1/4=3/4$. With
$d_\theta=1/5$, $c_\theta=2(1/2)(1/5)+1/25=0.24\le0.375=(1-\rho)(3/4)$.
Also $F_q(7/10)=\tfrac14(1.49)+\tfrac34(1.09)=1.19\le1.25-\tfrac12\cdot\tfrac25
\cdot\tfrac14=1.20$, so the first trial passes. The new costs are $1.49>h$ at
$\theta=0$ and $1.09\le h$ at $\theta=1$, so $C_J(K^+)=1/2$. The new margins are
$1-0.49=0.51$ and $1-0.09=0.91$.
\end{proof}

\subsection{Draft Stage~3 starting point: envelope accounting}
\label{subsec:stage3_envelope}

This proposition is a draft awaiting author review. Fix $K$, a direction $D\ne0$, and
$\bar s=\kappa/\|D\|_2$. Keep $\overline M_\theta$ from
Lemma~\ref{lem:stage2_cost_sensitivity}, and set $\beta_\theta=\|B(\theta)D\|_2$,
$\bar k=\|K\|_2+\kappa$, $w_0=\|Q\|_2+\|R\|_2\bar k^2$,
$w_1=2\|R\|_2\bar k\|D\|_2$, and $w_2=2\|R\|_2\|D\|_2^2$. For $z=(\theta,x)$, let
$X_0=\|x\|_2$, $U_0=V_0=0$, and
\[
 X_{t+1}=\overline M_\theta X_t,\quad
 U_{t+1}=\overline M_\theta U_t+\beta_\theta X_t,\quad
 V_{t+1}=\overline M_\theta V_t+2\beta_\theta U_t,
\]
\[
 C_z=\sum_{t=0}^{H-1}\big(2w_0(U_t^2+X_tV_t)+4w_1X_tU_t+w_2X_t^2\big),
\]
and define the \emph{alignment} and \emph{upper envelope}
\[
 a_z=-\langle\nabla_KJ_{H,K}(z),D\rangle_F,
 \qquad \mathcal U_z(s)=J_{H,K}(z)-sa_z+\tfrac12C_zs^2.
\]

\begin{proposition}[Envelope coverage bound]
\label{prop:stage3_envelope}
For every $s\in[0,\bar s]$, $J_{H,K+sD}(z)\le\mathcal U_z(s)$. Define the
\emph{guaranteed gains} and \emph{possible losses}
\[
 \mathcal A(s)=\{z\notin S_K^J:\mathcal U_z(s)\le h\},
 \qquad
 \mathcal R(s)=\{z\in S_K^J:\mathcal U_z(s)>h\}.
\]
Then
\[
 C_J(K+sD)-C_J(K)\ge\mu(\mathcal A(s))-\mu(\mathcal R(s)).
\]
Suppose $C_z>0$. A rejected $z$ with excess $e_z=J_{H,K}(z)-h>0$ is in
$\mathcal A(s)$ if and only if $a_z>0$, $a_z^2\ge2C_ze_z$, and
$s\in[(a_z-\sqrt{a_z^2-2C_ze_z})/C_z,(a_z+\sqrt{a_z^2-2C_ze_z})/C_z]$. An accepted
$z$ with slack $\sigma_z=h-J_{H,K}(z)\ge0$ is outside $\mathcal R(s)$ if and only if
$s\le(a_z+\sqrt{a_z^2+2C_z\sigma_z})/C_z$.
\end{proposition}

\begin{proof}
Let $r\in[0,\bar s]$, $K_r=K+rD$, and $G_r=G_K(\theta)-rB(\theta)D$, so
$\|G_r\|_2\le\overline M_\theta$ and $\|K_r\|_2\le\bar k$. Let $x_t=G_r^tx$, with
derivatives in $r$ marked by primes. Then
$x_{t+1}'=G_rx_t'-B(\theta)Dx_t$ and
$x_{t+1}''=G_rx_t''-2B(\theta)Dx_t'$, so induction gives $\|x_t\|_2\le X_t$,
$\|x_t'\|_2\le U_t$, and $\|x_t''\|_2\le V_t$. For $W_r=Q+K_r^\top RK_r$ we have
$\|W_r\|_2\le w_0$, $\|W_r'\|_2\le w_1$, and $\|W_r''\|_2\le w_2$. Since
\[
 (x_t^\top W_rx_t)''=2x_t'^\top W_rx_t'+2x_t^\top W_rx_t''
 +4x_t^\top W_r'x_t'+x_t^\top W_r''x_t,
\]
$\phi(r)=J_{H,K_r}(z)$ satisfies $|\phi''(r)|\le C_z$. Also $\phi'(0)=-a_z$, and
Taylor's formula with integral remainder gives
$\phi(s)\le\mathcal U_z(s)$. Thus $\mathcal A(s)$ is contained in the gained set, and
the lost set is contained in $\mathcal R(s)$. The bound follows
from~\eqref{eq:stage2_coverage_accounting}'s identity. The last two claims solve
$\tfrac12C_zs^2-a_zs+e_z\le0$ and $\tfrac12C_zs^2-a_zs-\sigma_z\le0$ over $s\ge0$.
For the first, the roots have positive product, so real roots are both positive
exactly when $a_z>0$.
\end{proof}

This bound uses old costs and gradients only, so it is computable before updated
costs are observed. Positivity is not automatic. On the 201 primary-family updates of
the recorded study it is never positive (largest value $-0.156$), although coverage
rises by up to $0.203$ on those updates. On the designed one-plant scalar fixture,
with $A=0.8$, $B=1$, $K=0.2$, $Q=1$, $R=0$, $H=2$, $h=1.2$, and four resets, the
boundary and high-cost targets give a bound of $0.25$, equal to the observed change,
while background gives $0$. This is a step-size effect, not evidence of target
superiority: with one plant every cost gradient has the same sign, and background's
shorter move $\Delta K\approx0.14$ leaves the largest cost at $1.21$, just above $h$,
whereas boundary's $\Delta K\approx0.22$ crosses it.

\begin{proposition}[Guard-limited updates depend only on direction]
\label{prop:stage3_direction}
Fix $K$ and a certificate $P$ with positive retained margins $\eta_\theta$, and let
\[
 \mathcal G=\{\Delta K:\|\Delta K\|_2\le\kappa,\
 c_\theta(\Delta K)\le(1-\rho)\eta_\theta
 \text{ for every }\theta\in\Theta_{\mathrm{ret}}\}.
\]
For $D\ne0$, let $s^*(D)=\max\{s\ge0:sD\in\mathcal G\}$, the \emph{step capacity}.
Then $s^*(D)$ is finite and positive, $\{s\ge0:sD\in\mathcal G\}=[0,s^*(D)]$, and for
every $\alpha>0$,
\[
 s^*(\alpha D)\,\alpha D=s^*(D)\,D.
\]
So the full-capacity update depends on $D$ only through $D/\|D\|_F$.
\end{proposition}

\begin{proof}
$\mathcal G$ is defined by conditions on $\Delta K$ alone. Along $s\mapsto sD$, the
norm $s\|D\|_2$ and each $c_\theta(sD)=2\|G_K(\theta)^\top P\|_2sd_\theta(D)
+\|P\|_2s^2d_\theta(D)^2$ are continuous and nondecreasing in $s\ge0$. They vanish at
$s=0$, where every budget is positive, and $s\|D\|_2$ exceeds $\kappa$ for large $s$.
Hence the feasible set is a compact interval $[0,s^*(D)]$ with $s^*(D)>0$. Finally,
$s\alpha D\in\mathcal G$ exactly when $s\alpha\le s^*(D)$, so
$s^*(\alpha D)=s^*(D)/\alpha$.
\end{proof}

The proposition ignores two limits on the step: the initial step $s_0$ and the Armijo
test. When either binds, the size of $D$ can matter. In the recorded family study the
$s_0$ limit almost never binds, but Armijo limits most updates on mirrored plants.
Under the proposition's conditions, a target can change the update only by turning
$D$, so any theorem that separates target rules needs a hypothesis on that turn.


\section{Numerical Results}
\label{sec:numerical_experiments}

\section{Discussion}
\label{sec:discussion}


\section{Conclusion}
\label{sec:conclusion}


\subsection*{Disclosure of AI Assistance}
\label{subsec:ai_disclosure}

\noindent Codex assisted with drafting and preliminary derivations.

\clearpage
\bibliography{references.bib}

\end{document}
